\subsection{Proper biset retracts and minimal models}

\begin{lemma}[A proper biset retract preserves a hypothetical projective plane]
\label{C:proper-biset-retract-CP2}
Let $S$ be a finite $(N,M)$-biset and let $f:S\to S$ be an
equivariant idempotent with nonempty image $U$.  If the collage
on $S$ has type $\mathbb CP^2$, then the collage on $U$ has
type $\mathbb CP^2$ as well.
\end{lemma}
\begin{proof}
The inclusion and $f$ give actual functors between the two
collages, identity on the vertices, whose composite on the
smaller collage is the identity.  Write their classifying
spaces as $X_U\subseteq X_S$ and their edge bars as
$E_U\subseteq E_S$.  These are inclusions of simplicial
subspaces.  The double mapping cylinder pair has quotient
\[
 X_S/X_U\simeq\Sigma(E_S/E_U).
\]
If restriction $H^2(X_S;\mathbb Q)\to H^2(X_U;\mathbb Q)$
were zero, its generator would lift to relative cohomology.
The suspension quotient has zero positive-degree products,
contradicting the projective-plane square.

Restriction is therefore nonzero.  The retraction splits
integral homology and the fundamental group.  Since those of
$X_S\simeq\mathbb CP^2$ are $\mathbb Z$ in degrees zero,
two and four, zero in the other positive degrees, and trivial
in degree one homotopy, the same holds for $X_U$ in degree two
and for its fundamental group.  Pullback from $X_U$ is an
isomorphism in degree two.  Cup naturality forces its fourth
homology to be nonzero, hence the full $\mathbb Z$ summand.
Thus the actual inclusion $X_U\to X_S$ is an integral
homology equivalence between simply connected CW spaces.
The homological Whitehead theorem makes it a homotopy
equivalence.
\end{proof}

\begin{corollary}\label{C:biset-core-prune}
A least-arrow two-object triangular model of $\mathbb CP^2$
has no proper equivariant biset retraction.  At sixteen arrows,
every proper equivariant idempotent on the cross-biset is
excluded by the already complete fifteen-arrow lower bound.
\end{corollary}

\begin{corollary}[A trivial action side gives an actual act-cone model]
\label{C:trivial-biset-side-CP2-retract}
If one side of a finite biset action is trivial and its collage
has type $\mathbb CP^2$, the corresponding ordinary act cone
also has type $\mathbb CP^2$.  A nontrivial trivially acting
corner can consequently be removed in a minimal model.
\end{corollary}
\begin{proof}
The collapse functor of Lemma~\ref{C:trivial-biset-side-H2}
has an actual section, obtained by inserting the identity of
the collapsed monoid.  Its pullback is an isomorphism on
rational $H^2$.  The split integral-homology and cup-product
argument in Lemma~\ref{C:proper-biset-retract-CP2} therefore
makes this section an integral homology equivalence between
simply connected CW spaces.
\end{proof}

\begin{lemma}[Commuting permutation actions]
\label{C:permutation-biset-cups}
If both actions on a nonempty finite biset are by permutations,
every rational degree-two product of its collage vanishes.
\end{lemma}
\begin{proof}
The edge action bar is a finite-sheeted covering of
$BM\times BN$.  Rational transfer makes pullback injective,
and K\"unneth embeds the direct sum of the two vertex
degree-two groups in the product's degree-two group.
Thus every collage degree-two class restricts to zero on
both vertices.  The relative suspension quotient kills all
their products, as in Theorem~\ref{C:two-cross-arrows-universal}.
\end{proof}
